\scanpage{25}
tels que $\forall j\geq1$, on ait comm. dans
\diagram{\begin{tikzcd}
M.T^{(i)} \arrow[r,".T^{(j)}"] \arrow[d,"u_i\otimes T^{(i)}"'] & MT^{(i+j)} \arrow[d,"u_{i+j}\otimes T^{(i+j)}"]\\
M'.T^{(i)} \arrow[r,".T^{(j)}"] & M'.T^{(i+j)}
\end{tikzcd}}
i.e. tels qu'on ait
\[
c_{ij}(u_{i+j}-u_i)=0\qquad\text{si }i\geq N,\ j\geq1,
\]
et le th. nous permet de conclure que tous les $u_i$ sont égaux, d'où la conclusion
\diagram{\begin{tikzcd}
\operatorname{Hom}_k(M,M') \arrow[r,"\sim"] & \operatorname{Hom}_{\unsure{S\text{-mod gr.d}|S\text{-mod gr.}}}(\omega_N(M),\omega_N(M'))
\end{tikzcd}}

\uline{Rmq 1} Le foncteur $\omega_N$ admet un adjoint à droite $\pi_N$
\[
\operatorname{Hom}_{\unsure{S\text{-mod gr.d}|S\text{-mod gr.}}}(\omega_N(M),\unsure{\mathcal M})\simeq\operatorname{Hom}_k(M,\pi_N(\mathcal M))
\]
en posant
\[
\begin{array}{l}
\pi_N(\mathcal M)\subset\displaystyle\prod_{i\geq N}\mathcal M_i\\
\qquad\text{''}\\
\displaystyle\left\{(x_i)_{i\geq N}\ \middle|\ x_iT^{(j)}=c_{ij}x_{i+j}\quad\forall i\geq N,\ j\geq1\right\}
\end{array}
\]
En considération de ceci, \deleted{il existe} les parties a) et b) du corollaire 1 sont équivalentes, b) signifie que l'hom. can. d'adjonction
\diagram{\begin{tikzcd}
M \arrow[r] & \pi_N(\omega_N(M))
\end{tikzcd}}
est \deleted{\unsure{tjs|tjrs}} un isomorphisme.

\uline{Rmq 2} Considérons $\omega_0(M)=M\otimes_k S=M\otimes_k k\{T\}$, alors $\omega_N(M)=S^{+(N)}.\omega_0(M)$. \deleted{Donc $\unsure{\omega_N}$ \unsure{en un sens}, $\unsure{\omega_0\to}$} \inserted{On peut aussi écrire, si $N\geq N'\geq0$}
\[
\omega_N(M)\simeq\deleted{\unsure{\prod}}\sum_{i\geq N}\omega_{N'}(M)_i\overset{\mathrm{def}}=\underset{\text{(``troncature'')}}{\tau_{\geq N}}(\omega_{N'}(M)).
\]
Donc en
\mspage{24}
